\title{Group Sequence Policy Optimization}

\begin{abstract}
We present Group Sequence Policy Optimization (GSPO), a reinforcement learning algorithm designed for large language models that excels in stability, efficiency, and overall effectiveness. Departing from traditional methods that rely on token-wise importance ratios, GSPO introduces an importance ratio calculated from the probability of the entire sequence, applying clipping, rewarding, and optimization at the sequence level. Experimental results reveal that GSPO outperforms the GRPO algorithm, delivers greater training efficiency, and enhances the reliability of Mixture-of-Experts (MoE) RL training, while also offering opportunities to streamline RL system architecture. These advantages of GSPO have played a key role in the notable advancements seen in the most recent Qwen3 models.
\end{abstract}

\section{Introduction}
\label{sec:intro}

Reinforcement learning (RL) has become a crucial framework for enhancing the scalability of language models \citep{o1,dpsk-r1,qwq32b,qwen3}. By leveraging RL at scale, these models acquire the proficiency to address complex tasks—including advanced mathematics and programming—via more profound and extended reasoning approaches.

A fundamental requirement for successfully increasing RL scale with higher computational resources is to preserve stable and resilient training dynamics. Nonetheless, leading RL algorithms such as GRPO \citep{grpo} are beset with significant stability problems, especially during the training of very large language models, which can frequently culminate in catastrophic and irreversible model failure \citep{qwen3, minimax-m1}. Such instability poses a significant challenge to further expanding the capabilities of language models through prolonged RL-based optimization.

This study reveals that the root cause of instability in GRPO is the improper implementation and subsequent breakdown of importance sampling weights within its algorithmic structure. As a result, significant high-variance noise is injected into the training process, intensifying as the length of generated responses grows, and further exacerbated by the application of clipping. This sequence of effects ultimately leads to catastrophic model failure.

To overcome these critical shortcomings, we introduce \textbf{Group Sequence Policy Optimization (GSPO)}, a novel reinforcement learning algorithm tailored for training large language models. The principal advancement of GSPO resides in its rigorously defined importance ratio—derived from sequence likelihood \citep{click}—thereby faithfully adhering to the foundational tenets of importance sampling. Moreover, GSPO leverages normalized rewards, calculated as the advantage across multiple responses for a given query, to ensure consistency between sequence-level reward evaluation and the optimization process.

Our experimental results clearly establish that GSPO considerably outperforms GRPO in terms of training stability, efficiency, and overall performance. Notably, GSPO effectively addresses the inherent instability issues present in reinforcement learning for large-scale Mixture-of-Experts (MoE) models, obviating the necessity for elaborate stabilization methods and potentially streamlining reinforcement learning infrastructure. These advantages directly facilitated the marked performance gains observed in the most recent Qwen3 models. We anticipate that GSPO will serve as a reliable and extensible algorithmic base, supporting ongoing progress in large-scale language model training using reinforcement learning.

\section{Preliminaries}

\paragraph{Notation}
Throughout this work, we represent an autoregressive language model with parameters $\theta$ as a policy $\pi_\theta$.
Let $x$ signify a query, while $\mathcal{D}$ refers to the collection of queries. For any query $x$ and its associated response $y$, the probability assigned to $y$ by the policy $\pi_\theta$ is written as $\pi_\theta (y | x)=\prod_{t=1}^{|y|} \pi_\theta (y_t | x, y_{<t} )$, where $|y|$ indicates the length of $y$ in tokens. Each query-response pair $(x, y)$ can be evaluated using a verifier function $r$, which yields a reward $r(x, y) \in [0, 1]$.

\paragraph{Proximal Policy Optimization (PPO)} Using samples generated from the old policy $\pi_{\theta_\text{old}}$, PPO \citep{ppo} constrains the policy update within a proximal region of the old policy through the clipping mechanism. Specifically, PPO employs the following objective for policy optimization (we omit the KL regularization term hereinafter for brevity, as it is not the focus of this paper):

\begin{align}
\mathcal{J}_\text{PPO}(\theta) = \mathbb{E}_{ x \sim \mathcal{D},\, y \sim \pi_{\theta_\text{old}}( \cdot | x) }
\left[ \frac{1}{|y|} \sum_{t=1}^{|y|} 
\min \left( w_{t}(\theta) \widehat{A}_{t},  \, \mathrm{clip} \left( w_{t}(\theta), 1 - {\varepsilon}, 1 + {\varepsilon}\right) \widehat{A}_{t} \right)
\right],
\end{align}

In this context, the importance ratio for token $y_t$ is expressed as $
w_{t}(\theta) = \frac{ \pi_{\theta} (y_{t} | x, y_{<t}) }{ \pi_{\theta_\text{old}} (y_{t} | x,y_{<t})} $. The advantage $\widehat{A}_{t}$ corresponding to $y_t$ is computed using an additional value model, while $\varepsilon$ represents the threshold for clipping the importance ratios.

A significant practical obstacle in using PPO is its pronounced dependence on the value model. Generally, the value model is constructed to be of comparable scale to the policy model, resulting in substantial demands for both memory and computation. The performance of the algorithm is closely tied to the accuracy of the value estimates. Attaining a dependable value model is already difficult, and scaling it to support longer generated responses or increasingly sophisticated tasks exacerbates these difficulties.

\paragraph{Group Relative Policy Optimization (GRPO)} GRPO \citep{grpo} bypasses the need for the value model by computing the relative advantage of each response within a group of responses to the same query. Specifically, GRPO optimizes the following objective:

\begin{align}
\label{equ:grpo}
\mathcal{J}_\text{GRPO}(\theta) = \mathbb{E}_{ x \sim \mathcal{D},\, \{y_i\}_{i=1}^G \sim \pi_{\theta_\text{old}}( \cdot | x) }
\left[ \frac{1}{G} \sum_{i=1}^{G} \frac{1}{|y_i|} \sum_{t=1}^{|y_i|} 
\min \left( w_{i,t}(\theta) \widehat{A}_{i,t},  \, \mathrm{clip} \left( w_{i,t}(\theta), 1 - {\varepsilon}, 1 + {\varepsilon}\right) \widehat{A}_{i,t} \right)
\right],
\end{align}

where $G$ is the number of generated responses to each query $x$ (i.e., the group size), and the importance ratio $w_{i,t}(\theta)$ and advantage $\widehat{A}_{i,t}$ of token $y_{i,t}$ are:

\begin{align}
    w_{i,t}(\theta)=\frac{ \pi_{\theta} (y_{i,t} | x, y_{i,<t}) }{ \pi_{\theta_\text{old}} (y_{i,t} | x,y_{i,<t})},\quad\ 
    \widehat{A}_{i,t} = \widehat{A}_{i} = \frac{r(x, y_i) - \mathrm{mean} \left( \{ r(x, y_i) \}_{i=1}^G \right) }{ \mathrm{std} \left( \{ r(x, y_i) \}_{i=1}^G \right) },
\end{align}

respectively, where all the tokens in $y_i$ share the same advantage as $\widehat{A}_{i}$.

\section{Motivation}
\label{sec:motivation}

The increasing complexity of models, including factors such as their overall size, the prevalence of sparsity (as seen in Mixture-of-Experts architectures), and the extension of response sequences, drives the need for substantial rollout batch sizes to ensure efficient hardware usage in RL applications. For the purpose of enhancing sample efficiency, it is a common approach to subdivide a large rollout batch into several smaller mini-batches for the purpose of performing gradient updates. However, this strategy unavoidably results in an off-policy learning scenario, given that responses $y$ are drawn according to a previous policy $\pi_{\theta_\text{old}}$, as opposed to the updated policy $\pi_\theta$ that is the subject of current optimization. This dynamic underlies the role of the clipping mechanism in algorithms such as PPO and GRPO, where it serves to exclude samples that have become excessively ``off-policy'' from contributing to the calculation of the gradient.

While mechanisms like clipping aim to manage this off-policy discrepancy, we identify a more fundamental issue in GRPO: \textit{its objective is ill-posed}. This problem becomes particularly acute when training large models on long-response tasks, leading to catastrophic model collapse. The ill-posed nature of the GRPO objective stems from a misapplication of importance sampling weights. The principle of importance sampling is to estimate the expectation of a function $f$ under a target distribution $\pi_\text{tar}$ by re-weighting samples drawn from a behavior distribution $\pi_\text{beh}$:

\begin{align}
\mathbb{E}_{z \sim \pi_\text{tar}} \left[ f(z) \right] = \mathbb{E}_{z \sim \pi_\text{beh}} \left[ \frac{ \pi_\text{tar} (z) }{ \pi_\text{beh} (z)} \, f(z) \right].
\end{align}

The effectiveness of correcting for distributional mismatch through the importance weight $\frac{ \pi_\text{tar} (z) }{ \pi_\text{beh} (z)}$ fundamentally depends on averaging across a large number of samples ($N \gg 1$) drawn from the behavior distribution $\pi_\text{beh}$. 

By contrast, GRPO employs the importance weight $\frac{ \pi_{\theta} (y_{i,t} | x, y_{i,<t}) }{ \pi_{\theta_\text{old}} (y_{i,t} | x,y_{i,<t})}$ at every individual token position $t$. Because this ratio is computed using just a single sampled token $y_{i,t}$ from each respective next-token distribution $\pi_{\theta_\text{old}}( \cdot | x, y_{i, <t})$, it does not achieve meaningful correction of the underlying distributional divergence. Rather, it results in substantial variance being injected into the training gradients, which tends to accumulate throughout lengthy sequences and is further intensified by the use of gradient clipping. Empirical evidence shows that this gradient noise can trigger a total collapse of the model, from which recovery proves to be nearly impossible; after collapse, attempts to restart training (including reloading prior checkpoints or carefully tuning hyperparameters like clipping ranges, increasing sequence length, or revising the RL queries) are generally ineffective.

This empirical finding points to an underlying flaw in the architecture of GRPO. The ineffectiveness of importance weighting at the token level underscores a critical conceptual issue: \textbf{the granularity of the optimization objective needs to correspond to the granularity of reward}. When rewards are distributed based on whole sequences, performing off-policy corrections at the token scale is not well-founded. This realization motivates a departure from token-level objectives and prompts a shift toward the use of importance weights and optimization performed on the \textit{sequence level}.

\section{Algorithm}

\subsection{GSPO: Group Sequence Policy Optimization}

Although the token-level importance weight $\frac{ \pi_{\theta} (y_{i,t} | x, y_{i,<t}) }{ \pi_{\theta_\text{old}} (y_{i,t} | x,y_{i,<t})}$ presents issues within GRPO, our analysis reveals that the \textit{sequence-level} importance weight $\frac{ \pi_{\theta} (y | x) }{ \pi_{\theta_\text{old}} (y | x)}$ holds significant theoretical relevance in language generation settings: it quantifies the extent to which the response $y$, drawn from $\pi_{\theta_\text{old}} (\cdot | x)$, differs from the distribution $\pi_{\theta} (\cdot | x)$. This measure not only directly corresponds to sequence-level rewards, but also functions as an effective metric for determining the application of the clipping mechanism.

Based on this straightforward observation, we propose the \textbf{Group Sequence Policy Optimization (GSPO)} algorithm. GSPO employs the following sequence-level optimization objective:

\begin{align}
\mathcal{J}_\text{GSPO} (\theta) =
\mathbb{E}_{ x \sim \mathcal{D},\, \{y_i\}_{i=1}^G \sim \pi_{\theta_\text{old}}( \cdot | x) }
\left[ 
\frac{1}{G} \sum_{i=1}^{G}
\min \left( s_{i}(\theta)  \widehat{A}_{i},  \, \mathrm{clip} \left( s_{i}(\theta), 1 - {\varepsilon}, 1 + {\varepsilon} \right) \widehat{A}_{i} \right) 
\right],
\label{equ:gspo}
\end{align}

where we adopt the group-based advantage estimation:

\begin{align}
\widehat{A}_{i} = \frac{r(x, y_i) - \mathrm{mean} \left( \{ r(x, y_i) \}_{i=1}^G \right) }{ \mathrm{std} \left( \{ r(x, y_i) \}_{i=1}^G \right) },
\end{align}

and define the importance ratio $s_{i}(\theta)$ based on sequence likelihood \citep{click}:

\begin{align}
s_{i}(\theta) = \left( \frac{ \pi_{\theta} (y_i | x) }{ \pi_{\theta_\text{old}} (y_i | x)} \right)^{\frac{1}{|y_i|}}
=
\exp \left( \frac{1}{|y_i|} \sum_{t=1}^{|y_i|} \log \frac{ \pi_{\theta} (y_{i,t} | x, y_{i,<t}) }{ \pi_{\theta_\text{old}} (y_{i,t} | x,y_{i,<t})} \right).
\end{align}

As a result, GSPO performs clipping at the level of complete responses, rather than individual tokens, effectively filtering out highly ``off-policy'' samples during gradient estimation. This strategy aligns with both the sequence-level reward structure and the approach to optimization. Additionally, we implement length normalization in $s_{i}(\theta)$ to stabilize variance and maintain $s_{i}(\theta)$ within a consistent numeric interval. Without such normalization, variations in a few token probabilities could lead to significant shifts in the sequence-level importance ratio, necessitating different clipping thresholds for responses based on their length. Furthermore, we observe that the clipping boundaries in GSPO, when compared with prior methods such as GRPO, generally vary substantially in magnitude, attributable to the different formulations of importance ratios in these techniques.

\subsection{Gradient Analysis}
\label{subsec:gradient}

We can derive the gradient of the GSPO objective as follows (clipping is omitted for brevity):

\begin{align}
\nabla_{\theta} \mathcal{J}_\text{GSPO} (\theta)
=&\ 
\nabla_{\theta} \mathbb{E}_{ x \sim \mathcal{D},\, \{y_i\}_{i=1}^G \sim \pi_{\theta_\text{old}}( \cdot | x) }
\left[ 
\frac{1}{G} \sum_{i=1}^{G} s_{i}(\theta) \widehat{A}_{i}
\right] \\
= &\ 
\mathbb{E}_{ x \sim \mathcal{D},\, \{y_i\}_{i=1}^G \sim \pi_{\theta_\text{old}}( \cdot | x) }
\left[ 
\frac{1}{G} \sum_{i=1}^{G}
s_{i}(\theta) \widehat{A}_{i}
\cdot \nabla_{\theta} \log s_{i}(\theta)
\right] \\
=&\ 
\mathbb{E}_{ x \sim \mathcal{D},\, \{y_i\}_{i=1}^G \sim \pi_{\theta_\text{old}}( \cdot | x) }
\left[ 
\frac{1}{G} \sum_{i=1}^{G} 
\left( \frac{ \pi_{\theta} (y_i | x) }{ \pi_{\theta_\text{old}} (y_i | x)} \right)^{\frac{1}{|y_i|}} \widehat{A}_{i} 
\cdot \frac{1}{|y_i|} \sum_{t=1}^{|y_i|} \nabla_{\theta} \log \pi_{\theta} (y_{i,t} | x, y_{i,<t}) 
\right].
\label{equ:gspo_grad}
\end{align}

For comparison, the gradient of the GRPO objective is as follows (note that $\widehat{A}_{i,t} = \widehat{A}_{i}$):

\begin{align}
\nabla_{\theta} \mathcal{J}_\text{GRPO} (\theta) 
=&\ 
\nabla_{\theta} \mathbb{E}_{ x \sim \mathcal{D},\, \{y_i\}_{i=1}^G \sim \pi_{\theta_\text{old}}( \cdot | x) }
\left[ \frac{1}{G} \sum_{i=1}^{G} \frac{1}{|y_i|} \sum_{t=1}^{|y_i|} 
w_{i,t}(\theta) \widehat{A}_{i,t}
\right] \\
=&\
\mathbb{E}_{ x \sim \mathcal{D},\, \{y_i\}_{i=1}^G \sim \pi_{\theta_\text{old}}( \cdot | x) }
\left[ \frac{1}{G} \sum_{i=1}^{G} \widehat{A}_{i} 
\cdot \frac{1}{|y_i|} \sum_{t=1}^{|y_i|} 
\frac{ \pi_{\theta} (y_{i,t} | x, y_{i,<t}) }{ \pi_{\theta_\text{old}} (y_{i,t} | x,y_{i,<t})} 
\nabla_{\theta} \log \pi_{\theta} (y_{i,t} | x, y_{i,<t})  
\right].
\end{align}

Consequently, the principal difference separating GSPO from GRPO concerns \textit{the method by which gradients of the log likelihoods for individual tokens are assigned weights}. GRPO employs a weighting scheme based on each token’s corresponding “importance weight” $\frac{ \pi_{\theta} (y_{i,t} | x, y_{i,<t}) }{ \pi_{\theta_\text{old}} (y_{i,t} | x,y_{i,<t})}$. The resultant weights are uneven, taking values in the interval $(0, 1+\varepsilon]$ (when $\widehat{A}_{i} >0$) or in $[1-\varepsilon, +\infty)$ (when $\widehat{A}_{i} < 0$). These disparities are substantial, and their cumulative effect throughout training may yield unforeseen outcomes. Conversely, GSPO treats all tokens within a generated response uniformly, thereby removing the instability present in GRPO’s approach.

\subsection{GSPO-token: A Token-level Objective Variant}

In scenarios like multi-turn RL, we may desire a finer-grained advantage adjustment than the sequence level. To this end, we introduce a token-level objective variant of GSPO, namely \textbf{GSPO-token}, to allow token-wise advantage customization:

\begin{align}
\mathcal{J}_\text{GSPO-token}(\theta)
=
\mathbb{E}_{ x \sim \mathcal{D},\, \{y_i\}_{i=1}^G \sim \pi_{\theta_\text{old}}( \cdot | x) }
\left[ 
\frac{1}{G} \sum_{i=1}^{G} \frac{1}{|y_i|} \sum_{t=1}^{|y_i|} 
\min \left( s_{i,t}(\theta) \widehat{A}_{i,t},  \, \mathrm{clip} \left( s_{i,t}(\theta), 1 - {\varepsilon}, 1 + {\varepsilon} \right) \widehat{A}_{i,t} \right) 
\right],
\label{equ:gspo-token}
\end{align}

where

\begin{align}
s_{i,t}(\theta) 
= \mathrm{sg} \left[ s_{i}(\theta) \right]  \cdot \frac{ \pi_{\theta} (y_{i,t} | x, y_{i,<t}) }{ \mathrm{sg} \left[ \pi_{\theta} (y_{i,t} | x, y_{i,<t}) \right] },
\end{align}

and $\mathrm{sg}[\cdot]$ denotes only taking the numerical value but stopping the gradient, corresponding to the \texttt{detach} operation in PyTorch. The gradient of GSPO-token can be derived as:

\begin{align}
\nabla_{\theta} \mathcal{J}_\text{GSPO-token}(\theta)
=&\ 
\nabla_{\theta} \mathbb{E}_{ x \sim \mathcal{D},\, \{y_i\}_{i=1}^G \sim \pi_{\theta_\text{old}}( \cdot | x) }
\left[ 
\frac{1}{G} \sum_{i=1}^{G}
\frac{1}{|y_i|} \sum_{t=1}^{|y_i|} 
s_{i,t}(\theta) \widehat{A}_{i,t}
\right] \\
=&\ 
\mathbb{E}_{ x \sim \mathcal{D},\, \{y_i\}_{i=1}^G \sim \pi_{\theta_\text{old}}( \cdot | x) }
\left[ 
\frac{1}{G} \sum_{i=1}^{G} s_i (\theta) \cdot \frac{1}{|y_i|} \sum_{t=1}^{|y_i|} 
\widehat{A}_{i,t} \frac{ \nabla_{\theta} \pi_{\theta} (y_{i,t} | x, y_{i,<t}) }{ \pi_{\theta} (y_{i,t} | x, y_{i,<t}) }
\right] \\
=&\ 
\mathbb{E}_{ x \sim \mathcal{D},\, \{y_i\}_{i=1}^G \sim \pi_{\theta_\text{old}}( \cdot | x) }
\left[ 
\frac{1}{G} \sum_{i=1}^{G}  \left( \frac{ \pi_{\theta} (y_i | x) }{ \pi_{\theta_\text{old}} (y_i | x)} \right)^{\frac{1}{|y_i|}}
\cdot \frac{1}{|y_i|} \sum_{t=1}^{|y_i|}
\widehat{A}_{i,t} \nabla_{\theta} \log \pi_{\theta} (y_{i,t} | x, y_{i,<t})  
\right].
\label{equ:gspo-token_grad}
\end{align}

It should be observed that the fraction $\frac{ \pi_{\theta} (y_{i,t} | x, y_{i,<t}) }{ \mathrm{sg} \left[ \pi_{\theta} (y_{i,t} | x, y_{i,<t}) \right] }$ evaluates to 1, thereby making $s_{i,t}(\theta)$ numerically equivalent to $s_{i}(\theta)$. By examining Equation~\eqref{equ:gspo} alongside Equation~\eqref{equ:gspo-token}, as well as Equation~\eqref{equ:gspo_grad} compared with Equation~\eqref{equ:gspo-token_grad}, it becomes clear that GSPO-token and GSPO yield the same values for the optimization target, the clipping criterion, and the corresponding theoretical gradient whenever the advantage for every token in $y_i$ is set uniformly (that is, $\widehat{A}_{i,t} = \widehat{A}_{i}$). Nevertheless, GSPO-token provides more adaptability, as it permits the assignment of distinct advantage values to each token.

\section{Experiments and Discussion}

\subsection{Empirical Results}

We conduct experiments using a cold-start model that is fine-tuned from Qwen3-30B-A3B-Base, and present both the training reward trajectories and performance metrics across the AIME'24 (mean Pass@1 out of 32 samples), LiveCodeBench (spanning 202410 to 202502, mean Pass@1 out of 8 samples), and CodeForces (Elo Rating) evaluation tasks. In the reinforcement learning (RL) process, each batch of rollout samples is split into four mini-batches for performing gradient updates. For GSPO, we configure the left and right clipping thresholds in Equation~\eqref{equ:gspo} as 3e-4 and 4e-4, correspondingly. As a baseline, we adopt GRPO, with its left and right clipping values in Equation~\eqref{equ:grpo} set to 0.2 and 0.27—values selected to facilitate an equitable comparison. It is important to highlight that GRPO requires the use of the Routing Replay training scheme for stable MoE RL convergence, a point elaborated in \S~\ref{subsec:routing_replay}, whereas \textit{GSPO eliminates this requirement}.

Figure~\ref{fig:results} illustrates that GSPO maintains stable progress during the entire training process. Notably, \textbf{GSPO facilitates ongoing enhancement of performance by scaling up training compute, periodically revising the query set, and increasing the length of generated outputs}. In addition, GSPO attains higher training efficiency compared to GRPO, evidenced by its improved accuracy and benchmark results when matched for training compute and query consumption. Ultimately, GSPO was effectively employed in RL training for the latest Qwen3 models, providing strong empirical support for GSPO's capability to harness the scaling benefits of RL for large language models.

\subsection{Curious Observation on Clipping Fractions}

One major difference between GSPO and GRPO lies in GSPO’s approach of clipping responses as wholes, as opposed to GRPO’s token-level clipping method. As illustrated in Figure~\ref{fig:clipping}, GSPO clips tokens at a rate approximately two orders of magnitude higher than GRPO; notably, this significant difference remains consistent even when the clipping ranges are modified. Surprisingly, despite GSPO utilizing far fewer tokens for training or for estimating gradients—due to the extensive clipping—it attains better training efficiency than GRPO. This unexpected result, whereby increased token clipping correlates with greater efficiency, suggests that the token-wise gradient estimations used by GRPO suffer from intrinsic noisiness and are not optimal for exploiting samples. Conversely, GSPO’s sequence-level methodology yields a more robust and informative signal for learning.

\subsection{Benefit of GSPO for MoE Training}
\label{subsec:routing_replay}

\paragraph{Background}
Relative to reinforcement learning (RL) training in standard dense models, Mixture of Experts (MoE) architectures exhibit particular instability due to their inherently sparse activation patterns. Specifically, when utilizing the GRPO algorithm, we observe that MoE models experience substantial \textit{expert-activation volatility}, which impedes the reliable convergence of RL training. Notably, even minor gradient updates may cause considerable shifts in which experts are selected to process identical inputs. To illustrate, in the case of the 48-layer Qwen3-30B-A3B-Base model, post each RL gradient step, approximately 10\% of the experts chosen under the revised policy $\pi_{\theta}$ differ from those under the previous policy $\pi_{\theta_\text{old}}$, for the same generated output. Such volatility, which intensifies with increased model depth, results in large fluctuations in the token-level importance weights $w_{i,t}(\theta) = \frac{ \pi_{\theta} (y_{i,t} | x, y_{i,<t}) }{ \pi_{\theta_\text{old}} (y_{i,t} | x,y_{i,<t})}$, rendering them unreliable as described in \S~\ref{sec:motivation} and~\ref{subsec:gradient}. This effect ultimately obstructs steady RL training convergence.

\paragraph{Our Previous Approach} In addressing this problem, our earlier method relied on the \textbf{Routing Replay} training scheme. This involved storing the activated experts generated by $\pi_{\theta_\text{old}}$ and then utilizing these cached routing patterns during the evaluation of $\pi_{\theta}$ while calculating the importance weights $w_{i,t}(\theta) = \frac{ \pi_{\theta} (y_{i,t} | x, y_{i,<t}) }{ \pi_{\theta_\text{old}} (y_{i,t} | x,y_{i,<t})}$. Consequently, for any token $y_{i,t}$, both $\pi_{\theta} (y_{i,t} | x, y_{i,<t})$ and $\pi_{\theta_\text{old}} (y_{i,t} | x,y_{i,<t})$ employ the identical set of network activations. This helps to maintain stable token-level importance ratios, thereby supporting consistent optimization of the active network configuration through successive gradient steps. As depicted in Figure~\ref{fig:routing_replay}, Routing Replay is a vital tool for achieving typical convergence behavior in GRPO training applied to MoE models.

\paragraph{Benefit of GSPO} While Routing Replay enables MoE models to train effectively under the GRPO framework, it does so at the cost of increased memory usage and communication, as well as the risk of artificially restricting the MoE model's capacity by forcing the reuse of routing modes. GSPO, by contrast, as illustrated in Figure~\ref{fig:results}, circumvents the necessity for Routing Replay. GSPO directly calculates importance ratios $s_i(\theta)$ in the standard manner, supports stable optimization, and achieves regular convergence. The crucial observation is that GSPO’s reliance is solely on the sequence-level likelihood, $\pi_{\theta} (y_i | x)$, rather than the likelihoods for individual tokens, $\pi_{\theta} (y_{i,t} | x, y_{i,<t})$, thereby insulating it from token-level variability. Since the MoE model consistently retains its language modeling prowess, the sequence likelihood remains relatively stable. In conclusion, GSPO directly tackles the issue of fluctuating expert activation in MoE models, eliminating the need for elaborate solutions such as Routing Replay. This facilitates a simpler, more robust training regime and allows the model to exploit its full capacity without imposed limitations.

\subsection{Benefit of GSPO for RL Infrastructure}

Owing to the differences in precision between engines utilized for training (such as Megatron) and those employed during inference (such as SGLang and vLLM), the standard approach involves recalculating sampled response likelihoods under the previous policy $\pi_{\theta_\text{old}}$ using the training engine. In contrast, GSPO's reliance on sequence-level likelihoods—rather than token-level likelihoods—renders it considerably more robust to such precision gaps. Consequently, GSPO enables direct usage of likelihood estimates produced by the inference engine within the optimization process, thus eliminating the necessity of recomputing them with the training engine. This property proves particularly advantageous in applications involving partial rollout, multi-turn RL, or when working within training-inference decoupled architectures.

\section{Conclusion}

We introduce Group Sequence Policy Optimization (GSPO), a novel reinforcement learning method tailored for large language model training. GSPO utilizes importance sampling principles, establishing importance ratios through sequence probability and applying sequence-wise operations such as clipping, rewarding, and optimization. Empirically, GSPO achieves significantly greater stability, training efficiency, and overall performance relative to GRPO, and proves especially beneficial in large-scale RL contexts for MoE architectures. This approach has underpinned major advancements observed in the latest Qwen3 models. By leveraging GSPO as a robust and scalable framework, we aim to further extend RL training at scale, anticipating substantial progress in artificial intelligence as a result.

\end{document}